\chapter{Fundamentos Teóricos}

Este capítulo nada mais é que um simples exemplo de como organizar um capítulo: não é obrigatório apresentar um arquivo chamado ``Fundamentos Teóricos''.

Note como é iniciado, primeiro uma breve explicação a seu respeito, e por último uma também breve descrição de seu conteúdo. Nas próximas seções sao apresentadas algumas dicas e recomendações sobre \LaTeX\ .

\section{\LaTeX}

O \LaTeX\ é um pacote de macros do \TeX\ que permite ao autor diagramar e imprimir seus trabalhos como documento de altíssima qualidade tipográfica, usando layout profissional.
O \LaTeX\ pode ser usado para produzir todos os tipos de documentos, desde uma simples carta até livros completos. O \LaTeX\ também é uma excelente ferramenta para desenvolver o \sigla{TC}{Trabalho de Curso}. 

\subsection{Extensões de arquivos \LaTeX}
O \LaTeX\ conta com diversos formatos de arquivo, os documentos são digitados dentro de um arquivo fonte com extensão .tex, que é processado gerando um arquivo .dvi, \sigla{DVI}{Device Independente File}. Outros formatos importantes:

\begin{itemize}

\item \textbf{.log, .aux:} guardam informações sobre a última vez que o LATEX processou seu documento, algumas vezes, modificações não aparecem após a compilação em função destes arquivos, pode ser necessário apagá-los para ter uma compilação inteiramente nova;

\item \textbf{.toc:} guarda os cabeçalhos. É usado para produzir o sumário;

\item \textbf{.lof:} É usado para produzir a lista de figuras;

\item \textbf{.lot:} É usado para produzir a lista de tabelas;

\item \textbf{.idx:} guarda todas as palavras que vão para o índice remissivo;

\item \textbf{.ind, .ilg:} usados para produzir o índice remissivo;

\item \textbf{.bib, .bbl:} usados para a bibliografia, .bib é a base bibliográfica onde são inseridas as bibliografias e .bbl é o arquivo gerado a partir do .bib, com as referências utilizadas no texto;

\item \textbf{.nlo:} usado para lista de siglas.

\end{itemize}

São gerados ainda arquivos no formato .pdf e .ps, sendo que a cada compilação a maioria destes arquivos são atualizados. 

\nota{Nunca apague arquivos com as extensões .tex, .bib e .sty deste modelo.}

\subsection{Aplicativos}

Diversas opções podem ser obtidas gratuitamente, sendo importante destacar a necessidade de duas partes em separado: os pacotes contendo as macros e o ambiente no qual os documentos são produzidos. 
Para a primeira parte, que também é instalada primeiro, recomenda-se o uso de: texlive(2013 ou superior) ou MikTeX(2.0.48 ou superior). Já para a segunda parte recomenda-se o uso de TexMaker(3.5.2 ou superior).

O uso de versões desatualizadas pode não gerar o resultado esperado ou mesmo não produzir resultado algum.



\subsection{Pacotes}

Os pacotes ou \textit{packages} são arquivos que implementam características adicionais para os documentos escritos em \LaTeX. Para inserir um pacote, digitamos:

 \verb+ \usepackage[opções]{pacote}+.

Abaixo estão listados os pacotes utilizados neste modelo:

\begin{itemize}
\item \textbf{inputenc} permite o uso de caracteres com til e cedilha;

\item \textbf{babel} A linguagem empregada nos títulos. Preenchendo o campo das opções com [brazil], todos os títulos são escritos em português. Se nada for feito, a linguagem usada é o inglês;

\item \textbf{fontenc} especifica a codificação;

\item \textbf{graphicx} pacote que permite a inserção de gráficos e figuras;

\item \textbf{geometry} pacote usado para definir as margens;

\item \textbf{hyperref} pacote que cria hiperlinks dentro do documento;

\item \textbf{fancyhdr} pacote usado para cabeçalhos e rodapés;

\item \textbf{mathptmx} fonte estilo times;

\item \textbf{abntex2cite}	usado para citação no formato abnt, pode ser usado com a opção [alf] ou [num], a primeira produz citações com alfanuméricos e a segunda citações através de números. No presente modelo a opção [alf] é utilizada, embora as duas sejam reconhecidas pela ABNT;

\item \textbf{formatacao} estilo particular deste \titulo, contempla cabeçalhos, formatação de títulos, ambientes como epígrafe e dedicatória;

\item \textbf{nomencl} utilizado para gerar a lista de acrônimos e siglas.

\item \textbf{indentfirst} identa a primeira linha de cada parágrafo.

\end{itemize}


Caso algum destes pacotes não esteja disponível no ambiente, ele deverá ser obtido automática ou manualmente, felizmente, a maioria dos editores buscam por pacotes que estejam faltando. A primeira compilação do modelo pode ser bastante demorada pois diversos pacotes serão buscados junto aos repositórios de macros \TeX.

No caso de não conseguir algum pacote de forma automática, é necessário baixá-lo e então proceder a instalação. No Windows, isso pode ser feito da seguinte maneira:


\begin{enumerate}
\item baixe o pacote;

\item cole tudo em arquivos de programas/miktex 2.9.48;

\item dê sim para todas as vezes que o sistema perguntar se deseja mesclar;

\item localize o programa miktex 2.9.48 e vá a maintenance(Admin) - Settings(Admin);

\item clique no botão Refresh FNDB.

\end{enumerate}



\subsection{Citação}

Monografia, assim como todo trabalho científico, deve ser fundamentada, deve conter citações que comprovem e garantam a validade do que está sendo dito. Contudo referenciar obras nem sempre é trivial, referências e citações devem ser feitas de maneira clara e consistente, por isso, este modelo de relatório utiliza ABN\TeX\ que é baseado nas normas: NBR6023:2002, \citeonline{NBR6023:2002} para referências, e NBR10520:2002, \citeonline{NBR10520:2002} para as citações. Observe alguns exemplos:

Segundo \citeonline{dubois}, formas podem ser classificadas....
 
Conforme \citeonline{rossatto}, `` O processo de taxonomia e identificação de plantas tem se tornado uma atividade sistemática'', \cite[p~1]{rossatto}.
 

Segundo\footnote{obviamente este exemplo é apenas ilustrativo} \apudonline[p~3]{dubois}{NBR6023:2002} as citações podem fazer uso de apud...
 \cite{smole}
É preciso observar cuidadosamente a forma de citar documentos \textit{online}, \cite{wikipedia}.

Trechos de citações direta devem seguir a norma ABNT que estabelece recuo de 4 centímetros, fonte 10, espaçamento simples no corpo do texto e dois espaços de 1,5 (tamanho 10) em cima e embaixo. Para facilitar utilize o comando\\ \verb+\citacaodireta{texto}{autor}{numero da pagina}+. 

\citacaodireta{A teleconferência permite ao indivíduo participar de um encontro nacional ou regional sem a necessidade de deixar seu local de origem. Tipos comuns de teleconferência incluem o uso da televisão, telefone e computador. Através de áudio-conferência, utilizando a companhia local de telefone, um sinal de áudio pode ser emitido em um salão de qualquer dimensão}{arica1}{1} 

\subsection{Inserção de Figuras}


Figuras devem ser referenciadas no texto, por exemplo, ``a figura \ref{logo} mostra o logo oficial do Instituto Federal Goiano...''. Para referenciar uma figura basta inserir um \textit{label} com o nome da figura dentro do ambiente figure e usar \verb+\ref{rotulodafigura}+ na hora de referenciar.

Figuras nunca devem ser referenciadas com as expressões abaixo, acima,  etc. O \LaTeX\ permite escolher entre topo, abaixo, centro, mas muitas vezes ele mesmo decide o melhor lugar. Por isso, coloque as figuras próximas ao local onde são referenciadas e utilize a forma citada anteriorente. Para incluir figuras utilize a seguinte forma: 

\begin{figure}[!htb] 
\begin{center}
\includegraphics[scale=1]{figuras/logo_oficial}
\caption{Logo oficial}\label{logo}
\end{center}
\end{figure}





Em que \textit{scale} é a escala desejada. São suportadas figuras em formato .jpeg, .eps. e .pdf. Para estes dois últimos casos pode ser necessário mais algum pacote.

\nota{Coloque todas as figuras na pasta figuras, assim o caminho para incluir as figuras será sempre o mesmo e tudo fica mais organizado.}

\subsection{Inserção de algoritmos}

Evite colocar trechos de códigos na monografia \footnote{Quando for realmente necessário coloque no anexo}, opte sempre por algoritmos, eles ficam mais elegantes e podem ser inseridos em qualquer parte do trabalho. Para facilitar a escrita de algoritmos utilize o pacote algorithm2e, cujo manual está incluído no modelo com o nome de ``manual\_algorithm2e''.

\verb+\usepackage[portuguese,algoruled,longend,linesnumbered]{algorithm2e}+. 

O resultado pode ser visto no algoritmo \ref{algoritmo};

%definindo os comandos Leia e Escreva pois não existem em portugues

\def\Leia{\textbf{leia }}
\def\Escreva{\textbf{escreva }}

\begin{algorithm}[hbtp] 
\label{algoritmo}
\caption{Ler número e imprimir se é par ou não.}
\Entrada{$numero$ - o número a ser verificado}
\Saida{Se o número é par ou não}

\Inicio{
\Leia($numero$)\;
\eSe{ $numero \% 2 = 0$ }
{\Escreva ($numero$, `` é par'')\;}
{\Escreva ($numero$, ``impar'')\;}
}

\end{algorithm}


\subsection{Notas de rodapés}

Notas de rodapé são uma excelente opção, elas esclarecem sem desviar o foco do leitor, por isso recomenda-se o uso delas. Mas cuidado, caso o texto fique excessivamente carregado de notas de rodapé, pode ser indício de que o trabalho não está bem organizado. Para adicionar uma nota de rodapé faça como no exemplo:

Em que f(x,y) é uma imagem suavizada\verb+\footnote{Através do filtro da média.}+

Produzindo o seguinte:

Em que f(x,y) é uma imagem suavizada\footnote{Através do filtro da média.}

As notas de rodapé são numeradas por capítulo, portanto, é de se esperar que haja mais de uma nota de rodapé de número 1.







\section{Considerações finais do capítulo}

Este capítulo apresentou conceitos fundamentais para a compreensão do trabalho. Dentre os principais, deve se destacar  ...

No próximo capítulo são abordados ... 


